\chapter{Colour}
As seen in the last "chapter", colour can be used as magnitude (value and saturation) and identity (hue) channels.
When using colour for nominal or ordinal data, it is important to have a limited palette which either is noncontiguous or a finite gradient.

For quantitative data using hue can go wrong, but may work.
When doing so, using the rainbow is a bad practice.
A smarter solution is to use sequential colormaps such as viridis, magma, or limit the gradient to two main hues.

\section{Colour Palettes}
There are different kinds of colour pallets to chose from.
\begin{description}
    \item[Categorical] Categorical colour palettes aim for maximum distinguishability between individual colours.

    \item[Univariate] Univariate colour palettes are for ordinal or quantitative data.
        They can be \emph{sequential}, \emph{diverging}, or \emph{cyclic}.
        Sequential colour palettes change either hue, saturation, or value.
        Diverging colour palettes change two dimensions of colour, one steadily increasing and one firstly increasing and then decreasing.
        Cyclic colour palettes increase and then decrease at least one colour dimension.

    \item[Bivariate] Bivariate colour palettes are univariate colour palettes which change at least two dimensions.
        A bivariate colour palette can e.g., change hue and saturation.
        Similarly to univariate colour palettes, bivariate colour palettes can also be sequential, diverging, or cyclic.

    \item[Binary] Binary colour palettes are special bivariate colour palettes, where there are only two values for on of the two dimensions.
\end{description}

\noindent
From these different kinds of palettes one must find the \emph{right} choice depending on the underlying data.

\subsection{Colour Vision Deficiency}
Colour vision deficiency (CVD) is the decreased ability to see colour or differences in colour, or distinguish shades of colour.
Colour palettes should be designed in mind of CVD, when many people see a graph.
Overall, at least 4.5\% of people of northern European decent have some kind of CVD.
Safe colour palettes to use with CVD in mind are those which only use blue and orange or yellow.
This maximises distinguishability for people with CVD.

\subsection{Colour Spaces}
There are different colour spaces to model colour perception.
Red-Green-Blue (RGB) is useful for displays but very unuseful for information encoding.
Instead often Hue-Saturation-Luminance (HSL) or Hue-Saturation-Value (HSV) is used to counteract this problem.
Cyan-Magenta-Yellow-Key (CMYK) is used for printing.

